% GENERATED FILE. Edit canonical lessons, then run npm run generate:books.
% canonical-source-hash: fnv1a64:5a9d13041a1864b6
% canonical-lessons: BN-C313-rajyo, BN-C313-somrat, BN-C313-singhason, BN-C313-mukut, BN-C313-shason-kora

\chapter{Kingdoms}
\label{ch:bn-kingdoms}
\hlchaptermodality{313}

\begin{chapteropening}
\textbf{By the end of this chapter:} \emph{I can talk about kingdoms, emperors and thrones.}

Complete the last lesson of chapter 313: I can talk about kingdoms, emperors and thrones.
\end{chapteropening}

\section[\texorpdfstring{\bn{রাজ্য} (rājyô)}{রাজ্য (rājyô)}]{\bn{রাজ্য} (rājyô) — a kingdom; a state}

\label{lesson:BN-C313-rajyo}



\begin{quote}
Before the new one: say the Bengali for ululation, then the Bengali for the groom's wedding party.
\end{quote}

\subsection*{You'll want to know: \bn{রাজ্য}}
\textbf{\bn{রাজ্য}} — \emph{rājyô} — \textquotedblleft{}a kingdom; a state\textquotedblright{}.

\begin{quote}
\bn{অনেক} \bn{দিন} \bn{আগে} \bn{এখানে} \bn{একটা} \bn{ছোট} \bn{রাজ্য} \bn{ছিল।}
\end{quote}

\begin{quote}
\emph{ônek din āge ekhāne ekṭā chhoṭo rājyô chhilô.}
\end{quote}

\begin{quote}
Long ago there was a small kingdom here.
\end{quote}

\begin{cousinweb}
\textbf{\bn{রাজ্য}} (\emph{rājyô}) is Sanskrit \emph{rājya}, 'kingship, sovereignty; a kingdom' (Monier-Williams), the abstract of \emph{rājan}, a king — \textbf{\bn{রাজা}}. Today it is also a state of India.
\end{cousinweb}

\begin{grammarlens}[title={What you've built}]
One more word for telling what happened.
\end{grammarlens}

\subsection*{Your turn}
\begin{itemize}
\raggedright
  \item \emph{Say it:} \emph{rājyô}
  \item \emph{Say it:} \emph{rājyô}, once more
  \item \emph{Say it:} \emph{ônek din āge ekhāne ekṭā chhoṭo rājyô chhilô.} — Long ago there was a small kingdom here.
  \item \emph{Recall:} say the Bengali for ululation, then the Bengali for the groom's wedding party, then say \emph{rājyô} again
\end{itemize}

\begin{tcolorbox}[breakable,colback=teal!4,colframe=teal!35!black,title={Before you move on}]
What does \bn{রাজ্য} mean? (\textquotedblleft{}A kingdom; a state\textquotedblright{}.) Say it once more.
\end{tcolorbox}
\section[\texorpdfstring{\bn{সম্রাট} (sômrāṭ)}{সম্রাট (sômrāṭ)}]{\bn{সম্রাট} (sômrāṭ) — an emperor}

\label{lesson:BN-C313-somrat}



\begin{quote}
Before the new one: say the Bengali for the groom's wedding party, then the Bengali for a kingdom.
\end{quote}

\subsection*{You'll want to know: \bn{সম্রাট}}
\textbf{\bn{সম্রাট}} — \emph{sômrāṭ} — \textquotedblleft{}an emperor\textquotedblright{}.

\begin{quote}
\bn{সম্রাট} \bn{দুর্গ} \bn{থেকে} \bn{পুরো} \bn{রাজ্য} \bn{দেখতেন।}
\end{quote}

\begin{quote}
\emph{sômrāṭ durgô theke puro rājyô dekhten.}
\end{quote}

\begin{quote}
The emperor used to look out over the whole kingdom from his fort.
\end{quote}

\begin{cousinweb}
\textbf{\bn{সম্রাট}} (\emph{sômrāṭ}) is Sanskrit \emph{samrāj}, 'a universal ruler, a paramount sovereign' (Monier-Williams), in its nominative \emph{samrāṭ}; an emperor rules over kings.
\end{cousinweb}

\begin{grammarlens}[title={What you've built}]
One more word for telling what happened.
\end{grammarlens}

\subsection*{Your turn}
\begin{itemize}
\raggedright
  \item \emph{Say it:} \emph{sômrāṭ}
  \item \emph{Say it:} \emph{sômrāṭ}, once more
  \item \emph{Say it:} \emph{sômrāṭ durgô theke puro rājyô dekhten.} — The emperor used to look out over the whole kingdom from his fort.
  \item \emph{Recall:} say the Bengali for the groom's wedding party, then the Bengali for a kingdom, then say \emph{sômrāṭ} again
\end{itemize}

\begin{tcolorbox}[breakable,colback=teal!4,colframe=teal!35!black,title={Before you move on}]
What does \bn{সম্রাট} mean? (\textquotedblleft{}An emperor\textquotedblright{}.) Say it once more.
\end{tcolorbox}
\section[\texorpdfstring{\bn{সিংহাসন} (singhāsôn)}{সিংহাসন (singhāsôn)}]{\bn{সিংহাসন} (singhāsôn) — a throne}

\label{lesson:BN-C313-singhason}



\begin{quote}
Before the new one: say the Bengali for a kingdom, then the Bengali for an emperor.
\end{quote}

\subsection*{You'll want to know: \bn{সিংহাসন}}
\textbf{\bn{সিংহাসন}} — \emph{singhāsôn} — \textquotedblleft{}a throne\textquotedblright{}.

\begin{quote}
\bn{নতুন} \bn{রাজা} \bn{সিংহাসনে} \bn{বসলেন।}
\end{quote}

\begin{quote}
\emph{nôtun rājā singhāsône bôslen.}
\end{quote}

\begin{quote}
The new king took the throne.
\end{quote}

\begin{cousinweb}
\textbf{\bn{সিংহাসন}} (\emph{singhāsôn}) is Sanskrit \emph{siṃhāsana}, 'a lion's seat, a king's seat' (Monier-Williams): \textbf{\bn{সিংহ}}, a lion, and \textbf{\bn{আসন}}, a seat.
\end{cousinweb}

\begin{grammarlens}[title={What you've built}]
One more word for telling what happened.
\end{grammarlens}

\subsection*{Your turn}
\begin{itemize}
\raggedright
  \item \emph{Say it:} \emph{singhāsôn}
  \item \emph{Say it:} \emph{singhāsôn}, once more
  \item \emph{Say it:} \emph{nôtun rājā singhāsône bôslen.} — The new king took the throne.
  \item \emph{Recall:} say the Bengali for a kingdom, then the Bengali for an emperor, then say \emph{singhāsôn} again
\end{itemize}

\begin{tcolorbox}[breakable,colback=teal!4,colframe=teal!35!black,title={Before you move on}]
What does \bn{সিংহাসন} mean? (\textquotedblleft{}A throne\textquotedblright{}.) Say it once more.
\end{tcolorbox}
\section[\texorpdfstring{\bn{মুকুট} (mukuṭ)}{মুকুট (mukuṭ)}]{\bn{মুকুট} (mukuṭ) — a crown}

\label{lesson:BN-C313-mukut}



\begin{quote}
Before the new one: say the Bengali for an emperor, then the Bengali for a throne.
\end{quote}

\subsection*{You'll want to know: \bn{মুকুট}}
\textbf{\bn{মুকুট}} — \emph{mukuṭ} — \textquotedblleft{}a crown\textquotedblright{}.

\begin{quote}
\bn{রানির} \bn{মুকুট} \bn{ঝলমল} \bn{করছিল।}
\end{quote}

\begin{quote}
\emph{rānir mukuṭ jhôlmôl kôrchhilô.}
\end{quote}

\begin{quote}
The queen's crown was glittering.
\end{quote}

\begin{cousinweb}
\textbf{\bn{মুকুট}} (\emph{mukuṭ}) is Sanskrit \emph{mukuṭa}, 'a tiara, a diadem, a crown' (Monier-Williams).
\end{cousinweb}

\begin{grammarlens}[title={What you've built}]
One more word for telling what happened.
\end{grammarlens}

\subsection*{Your turn}
\begin{itemize}
\raggedright
  \item \emph{Say it:} \emph{mukuṭ}
  \item \emph{Say it:} \emph{mukuṭ}, once more
  \item \emph{Say it:} \emph{rānir mukuṭ jhôlmôl kôrchhilô.} — The queen's crown was glittering.
  \item \emph{Recall:} say the Bengali for an emperor, then the Bengali for a throne, then say \emph{mukuṭ} again
\end{itemize}

\begin{tcolorbox}[breakable,colback=teal!4,colframe=teal!35!black,title={Before you move on}]
What does \bn{মুকুট} mean? (\textquotedblleft{}A crown\textquotedblright{}.) Say it once more.
\end{tcolorbox}
\section[\texorpdfstring{\bn{শাসন} \bn{করা} (shāsôn kôrā)}{শাসন করা (shāsôn kôrā)}]{\bn{শাসন} \bn{করা} (shāsôn kôrā) — to rule; to govern}

\label{lesson:BN-C313-shason-kora}



\begin{quote}
Before the new one: say the Bengali for a throne, then the Bengali for a crown.
\end{quote}

\subsection*{You'll want to know: \bn{শাসন} \bn{করা}}
\textbf{\bn{শাসন} \bn{করা}} — \emph{shāsôn kôrā} — \textquotedblleft{}to rule; to govern\textquotedblright{}.

\begin{quote}
\bn{সেই} \bn{রাজা} \bn{চল্লিশ} \bn{বছর} \bn{রাজ্য} \bn{শাসন} \bn{করলেন।}
\end{quote}

\begin{quote}
\emph{sei rājā chôllish bôchhôr rājyô shāsôn kôrlen.}
\end{quote}

\begin{quote}
That king ruled the kingdom for forty years.
\end{quote}

Past (he/she …): \bn{তিনি} \bn{দেশ} \bn{শাসন} \bn{করলেন}, \emph{tini desh shāsôn kôrlen}.

\begin{cousinweb}
\textbf{\bn{শাসন} \bn{করা}} (\emph{shāsôn kôrā}) puts Sanskrit \emph{śāsana}, 'chastising; teaching; an order' (Monier-Williams), before \bn{করা}. It shares its root with \textbf{\bn{শাসানো}}, to threaten, and \textbf{\bn{শাস্তি}}, a punishment.
\end{cousinweb}

\begin{grammarlens}[title={What you've built}]
One more word for telling what happened.
\end{grammarlens}

\subsection*{Your turn}
\begin{itemize}
\raggedright
  \item \emph{Say it:} \emph{shāsôn kôrā}
  \item \emph{Say it:} \emph{shāsôn kôrā}, once more
  \item \emph{Say it:} \emph{sei rājā chôllish bôchhôr rājyô shāsôn kôrlen.} — That king ruled the kingdom for forty years.
  \item \emph{Recall:} say the Bengali for a throne, then the Bengali for a crown, then say \emph{shāsôn kôrā} again
\end{itemize}

\begin{tcolorbox}[breakable,colback=teal!4,colframe=teal!35!black,title={Before you move on}]
What does \bn{শাসন} \bn{করা} mean? (\textquotedblleft{}To rule; to govern\textquotedblright{}.) Say it once more.
\end{tcolorbox}
